\section{Results of the Extended Evaluation}\label{sec:extended_results}
The extension recovers functionality on many previously problematic
payloads. In the selected recovery cohort, the x64 working rate rises
from 1.2\% to 80.2\%; the corresponding x86 rates are 14.9\% and 25.9\%
(Section~\ref{sec:recovery}). The following experiments also measure
data identification, controlled transformation behavior, and the new wrapping
and LLM modes on the full original corpus. Each experiment answers a
different question and retains its own input set.

\subsection{Data Identification: Known Answers and Original Payloads}
On the six classification fixtures, all 18 LLM invocations and all
six symbolic runs return the exact reference interval sets
(Table~\ref{tab:new_classifier}). Both methods miss zero data bytes and
classify zero code bytes as data. The three fixtures containing data have
precision and recall equal to one; the other three correctly return no
data. All resulting executions pass their buffer checks. This includes
the deliberately misleading printable-code case, which is rejected as
data rather than extracted as a string.


\begin{table}[t]
\caption{Six programs with independently annotated data locations. $D$ is annotated data bytes.
Exact interval sets are compared against predeclared references.
Mean times are per run in seconds.}
\label{tab:new_classifier}
\centering\small\setlength{\tabcolsep}{4pt}
\begin{tabular}{@{}llrrrrr@{}}
\toprule
Arch. & Case & $D$ & LLM exact & Symbolic exact & LLM time & Sym. time \\
\midrule
x64 & databait & 0 & 3/3 & 1/1 & 4.6 & 2.4 \\
x64 & multidata & 11 & 3/3 & 1/1 & 4.2 & 2.6 \\
x64 & nodata & 0 & 3/3 & 1/1 & 2.1 & 2.4 \\
x64 & table & 16 & 3/3 & 1/1 & 2.6 & 2.4 \\
x86 & multidata & 11 & 3/3 & 1/1 & 4.4 & 2.5 \\
x86 & nodata & 0 & 3/3 & 1/1 & 1.9 & 2.3 \\
\bottomrule
\end{tabular}
\end{table}


Both methods correctly identify all six controlled examples. The
comparison on 19 original payloads then evaluates how the proposed
regions are handled in the complete pipeline. Both pipelines build 19/19 outputs and pass
17/19 harness checks. Their compared intervals agree in 9/19 cases,
while their final binaries are byte-identical in 16/19. Shared filtering
and reference handling can therefore remove differences between proposals.
These two values describe agreement between the methods, whereas
accuracy against known answers is measured on the annotated programs. The same two inputs fail in both arms, which is consistent with
a shared downstream limitation but does not by itself isolate its cause.

The recorded median durations are 123.95 seconds for symbolic preparation
and 59.88 seconds for the LLM arm. The latter's estimated cost is \$0.2218
over 19 runs. These are measured run durations; preparation and caching differ
between the two arms. On the separate six-fixture benchmark,
the total LLM cost is approximately \$0.0059. 

\subsection{Functional Tests on 13 Controlled Programs}
The functional benchmark asks whether a transformed program still produces
its expected output. Each native configuration is run three times per
input: 21 tests on the 7 x64 programs and 18 on the 6 x86 programs.
Table~\ref{tab:controlled_functional} shows successful tests over planned
tests. It combines the two LLM prompting conditions to keep the functional
comparison readable; their separate results are supplied in the supplement.

\begin{table}[t]
\caption{Functional success on the controlled programs. Each entry is
successful tests/planned tests. LLM columns combine zero-shot and two-shot
requests at temperature 0.2. PII is outside this benchmark; wrapping is
native-only. The native x86 IBP row refers to the recorded earlier version.}
\label{tab:controlled_functional}
\centering\small\setlength{\tabcolsep}{7pt}
\begin{tabular}{@{}lrrrr@{}}
\toprule
 & \multicolumn{2}{c}{Native} & \multicolumn{2}{c}{LLM} \\
\cmidrule(lr){2-3}\cmidrule(lr){4-5}
Family & x64 & x86 & x64 & x86 \\
\midrule
EQ & 21/21 & 18/18 & 42/42 & 36/36 \\
NOP & 21/21 & 18/18 & 42/42 & 36/36 \\
BCF & 21/21 & 18/18 & 42/42 & 36/36 \\
IBP & 21/21 & 0/18 & 41/42 & 30/36 \\
Wrapping & 21/21 & 18/18 & -- & -- \\
\bottomrule
\end{tabular}
\end{table}

All native x64 tests pass. Native x86 EQ, NOP, BCF, and wrapping also
pass; the recorded native x86 IBP configuration produces 17 binaries from
18 attempts, but none passes the complete-buffer test.

The LLM produces 305 accepted outputs from 312 requests (97.8\%), and
every accepted output has the correct buffer. The remaining seven requests
comprise five structural-validation rejections and two transport errors.
Thus the main result is functional preservation for all accepted outputs
on these controlled inputs, with construction availability recorded separately.

Repeated output hashes are most visible in EQ. These inputs are short,
offer few locations for substitutions, and are subject to the restricted
set of supported transformations. This provides a straightforward
explanation for limited diversity on this suite; different requests need
not produce different accepted instruction sequences. Detailed hashes,
per-shot outcomes, time, and cost remain available in the supplement.

\subsection{Wrapping on All 102 Original Payloads}
The full-corpus campaign attempts 306 variants, three per input. It constructs
243 binaries (79.4\% of planned cases), of which 170 pass the harness.
Thus $W/N=55.6$\%, whereas $W/B=70.0$\%. These rates report both availability of a transformed binary and
its functional outcome.


\begin{table}[t]
\caption{Wrapping on all 102 original payloads, with three seeds per input.
$U$ counts distinct constructed hashes.}
\label{tab:new_wrap_full}
\centering\small
\begin{tabular}{@{}lrrrrr@{}}
\toprule
Arch. & Inputs & $B/N$ & $W/N$ & $W/B$ & $U$ \\
\midrule
x64 & 38 & 63/114 & 63/114 & 63/63 & 63 \\
x86 & 64 & 180/192 & 107/192 & 107/180 & 144 \\
\bottomrule
\end{tabular}
\end{table}


The two architectures fail at different stages. On x64, 51 attempts do
not build, while all 63 constructed outputs pass. On x86, only twelve
attempts do not build, but 73 of the 180 constructed outputs fail the
harness. All 73 also fail the saved isolated retests; those retests are
not counted as additional variants. There are 207 distinct hashes among
the 243 binaries, showing that different seeds do not always produce
different bytes.

The construction logs include incomplete instruction interpretation and
inconsistent address references. These failure stages have a direct
relationship to preparation: interpreting embedded data as instructions
can prevent a complete valid disassembly, and a reference that no longer
points to the intended location fails the consistency check after layout
changes. The logs identify these symptoms; they do not assign one common
cause to every failed execution.

\subsection{Completed LLM Transformation on the Original Corpus}
All 816 planned cells have an outcome. The campaign constructs 350
outputs, and 275 pass the external harness: a build rate of 42.9\%, an
overall working rate of 33.7\%, and a working rate among builds of 78.6\%.
Table~\ref{tab:payload_llm} shows the complete architecture/family/shot
breakdown. These are pipeline outcomes, including the retries described
in Section~\ref{sec:methodology}.


\begin{table}[t]
\caption{LLM transformation of all 102 original payloads. $N$: planned cells; $B$: accepted
builds; $W$: builds passing the external harness; $U$: distinct built-output
hashes within the condition. Every planned cell is retained.}
\label{tab:payload_llm}
\centering\small
\begin{tabular}{@{}lllrrrr@{}}
\toprule
Arch. & Family & Shots & $N$ & $B$ & $W$ & $U$ \\
\midrule
x64 & EQ & 0 & 38 & 7 & 7 & 7 \\
x64 & EQ & 2 & 38 & 6 & 6 & 6 \\
x64 & NOP & 0 & 38 & 8 & 8 & 8 \\
x64 & NOP & 2 & 38 & 6 & 6 & 6 \\
x64 & BCF & 0 & 38 & 6 & 6 & 6 \\
x64 & BCF & 2 & 38 & 6 & 5 & 6 \\
x64 & IBP & 0 & 38 & 6 & 6 & 6 \\
x64 & IBP & 2 & 38 & 8 & 8 & 8 \\
\midrule
x86 & EQ & 0 & 64 & 34 & 25 & 30 \\
x86 & EQ & 2 & 64 & 36 & 28 & 34 \\
x86 & NOP & 0 & 64 & 38 & 28 & 38 \\
x86 & NOP & 2 & 64 & 35 & 27 & 35 \\
x86 & BCF & 0 & 64 & 40 & 31 & 40 \\
x86 & BCF & 2 & 64 & 36 & 26 & 36 \\
x86 & IBP & 0 & 64 & 39 & 28 & 36 \\
x86 & IBP & 2 & 64 & 39 & 30 & 38 \\
\bottomrule
\end{tabular}
\end{table}



On x64, only 53/304 cells produce binaries, but 52/53 of those binaries
pass. On x86, 297/512 produce binaries and 223/297 pass. Construction is
therefore the principal observed bottleneck on x64, while x86 also has a
substantial failure rate after construction. The controlled programs isolate small instruction and data patterns;
the original payloads add larger listings and more complex layouts.
The two experiments therefore measure different levels of difficulty.

Zero-shot produces 178 builds and 139 passing outputs from 408 cells;
two-shot produces 172 builds and 136 passing outputs from 408 cells.
Their overall working rates are 34.1\% and 33.3\%. These descriptive
differences offer no consistent evidence that adding two examples improves
the outcome. They also concern repeated configurations of the same inputs,
not 816 independent programs.

The 466 failed cells have recorded error categories comprising 245
disassembly/boundary errors, 74 validation errors, 53 output-length
terminations, one transport error, and 93 other errors. These labels
describe the observed failure stage, not proven root causes. In particular,
truncation means the response reaches its output limit before delivering
a complete acceptable listing; it is distinct from a generated program
failing at execution time.

The 350 built outputs contain 323 distinct hashes. The median retained-run duration is
160.3 seconds, including local processing. Summing the representative
record for each cell gives an estimated \$19.2700, but including the
14 restarted failed runs raises recorded campaign cost to
\$19.6258. The latter is the appropriate total for work actually recorded;
missing usage for one run prevents claiming complete billing coverage.

\subsection{Updated Working Rates by Transformation}
\label{sec:recovery}
The recovery campaign retests 49 previously low-rate inputs: 13 x64 and
36 x86. On this selected subset, the x64 working rate rises from
23/1,941 (1.2\%) to 1,564/1,950 (80.2\%). The x86 rate rises from
772/5,195 (14.9\%) to 699/2,700 (25.9\%). These observations measure
the updated implementation and test harness together.

Table~\ref{tab:overall_working_update} places the recovery results back
into the overall historical population. For each transformation, the
retested inputs use their updated rates; other inputs retain their
historical rates. We keep the original input weights so that the reduction
from generally ten to five x86 repetitions cannot increase the overall
rate merely by reducing the influence of difficult inputs.

The reconstructed mean across the 15 configurations increases from
55.6\% to 82.6\% on x64 and from 45.0\% to 50.3\% on x86.
These are summaries of combined historical and updated evidence. The
original published figure is retained as the baseline; the update appears
here with its construction stated explicitly.

\begin{table}[t]
\caption{Overall working rates (\%): historical measurements and a
reconstructed update using the recovery tests. The update replaces rates
only for the 49 retested inputs and retains historical rates for the rest.
Historical input weights are held fixed. The final row is the unweighted
mean over the 15 configurations, as in the historical presentation.
This is a combined evidence summary, not a new execution of the full corpus.}
\label{tab:overall_working_update}
\centering\small\setlength{\tabcolsep}{5pt}
\begin{tabular}{@{}lrrrr@{}}
\toprule
 & \multicolumn{2}{c}{x64} & \multicolumn{2}{c}{x86} \\
\cmidrule(lr){2-3}\cmidrule(lr){4-5}
Configuration & Historical & Updated & Historical & Updated \\
\midrule
NOP & 65.5 & 92.1 & 48.3 & 49.1 \\
EQ & 65.5 & 99.7 & 77.5 & 76.2 \\
PII & 60.3 & 89.2 & 38.8 & 43.9 \\
BCF & 65.8 & 94.5 & 52.7 & 66.1 \\
IBP & 65.8 & 94.2 & 51.7 & 57.2 \\
EQ+NOP+PII & 64.2 & 86.3 & 24.0 & 26.9 \\
EQ+NOP+BCF & 1.1 & 29.5 & 43.9 & 54.7 \\
EQ+NOP+IBP & 65.8 & 91.3 & 40.8 & 42.8 \\
NOP+EQ & 66.0 & 90.5 & 48.1 & 49.2 \\
NOP+PII & 57.3 & 80.1 & 29.3 & 32.5 \\
NOP+BCF & 0.3 & 28.8 & 45.8 & 53.3 \\
NOP+IBP & 66.3 & 89.9 & 40.5 & 44.7 \\
EQ+PII & 57.3 & 86.5 & 28.0 & 31.7 \\
EQ+BCF & 66.8 & 92.6 & 55.2 & 67.3 \\
EQ+IBP & 66.0 & 94.4 & 50.6 & 58.4 \\
\midrule
\textbf{Mean} & 55.6 & 82.6 & 45.0 & 50.3 \\
\bottomrule
\end{tabular}
\end{table}


The lower x86 recovery rate concerns a deliberately difficult subset,
not all x86 inputs. Of its 36 inputs, 13 still have no successful
transformed variant. Preparation-only controls pass 340/360 tests (94.4\%),
whereas adding the native transformations yields 699/2,700 (25.9\%).
This separates successful preparation from the additional transformation
failures; it does not support attributing every remaining failure to data
identification. The dominant recorded execution symptom is an access
violation, which alone does not identify the underlying implementation error.

The improvement is substantially stronger on x64. Even excluding
MessageBox, whose test recognition was corrected, the selected x64 rate
rises from 19/1,791 (1.1\%) to 1,414/1,800 (78.6\%). Per-input and
per-configuration counts, including the selected-subset comparison, remain
available in the supplementary records.

\subsection{Overall Detection Rates for Wrapping and LLM Transformation}
The detection evaluation covers the outputs from the full 102-input LLM
campaign, not only the separate 19-input data-identification experiment.
Of the 816 planned LLM cells, 350 produce binaries, representing 323
distinct hashes. Table~\ref{tab:current_scans} reports the mean engine
detection rate using the same definition as the historical paper.

\begin{table}[t]
\caption{Static detection of the current outputs. Files: produced binaries;
hashes: distinct contents; reports: nonempty scan results. DR is the mean
percentage of engine responses marked malicious.
The working-output row is a subset of the overall LLM row. Report
completeness and response-category details are supplied in the supplement.}
\label{tab:current_scans}
\centering\small\setlength{\tabcolsep}{6pt}
\begin{tabular}{@{}lrrrr@{}}
\toprule
Condition & Files & Hashes & Reports & DR (\%) \\
\midrule
Wrapping & 243 & 207 & 207 & 0.00 \\
LLM x64 & 53 & 51 & 51 & 8.86 \\
LLM x86 & 297 & 272 & 271 & 4.35 \\
\midrule
LLM overall & 350 & 323 & 322 & 5.06 \\
LLM working outputs & 275 & 263 & 263 & 5.89 \\
\bottomrule
\end{tabular}
\end{table}

The overall LLM mean detection rate is 5.06\% across 322 nonempty reports,
equivalent to 3.80 malicious engine responses per file. The 275 outputs
that pass the functional test correspond to 263 distinct hashes and have
a mean detection rate of 5.89\%. Showing both rows distinguishes all
constructed outputs from the functionally successful subset.

All 207 distinct wrapping outputs have nonempty reports with no malicious
responses. These are static raw-file measurements; the engine-response
denominator follows the historical convention described in the methodology.
The separate data-identification scans and earlier pilot campaigns are
archived in the supplementary evidence rather than pooled into these results.
